\section{Related work}
\label{sec:related}
\subsection{Calibrated supervised and benign-deviation evidence}
A close comparison is the Mode-Calibrated Dual-Evidence Detector (MCDE) of Yue et al.~\cite{yue2026mcde}. It combines a supervised histogram-gradient-boosting classifier with distances to learned benign modes, converts both scores to empirical benign-tail evidence, allocates branch shares, and calibrates a final fused threshold on a separate benign set. Its principal evaluation withholds entire attack families. A supplementary device-transfer experiment recalibrates the final threshold with trusted target benign data and exposes a false-alert versus recall tradeoff.

MCDE establishes precedent for calibrated supervised and benign-deviation fusion. PACE retains multiple foreign peer classifiers as separately thresholded experts and adds a heterogeneous local anchor and bank. It calibrates the peer union, anchor tier, and bank union separately, without a final aggregate threshold, and its low-alpha resolution rule does not preserve a strict partition of the nominal allocation. Its target-label restriction also differs from family hold-out: peers may have seen the target's attack families. The resulting comparison concerns the retained experts, the final decision structure, and the target information regime; MCDE was not run in the PACE campaign.

Table~\ref{tab:prior} summarizes the closest architectural and information-regime distinctions.
\begin{table}[tbp]
\centering
\small
\caption{Methodological comparison of information requirements and decision mechanisms, not a performance ranking. None of the external systems was rerun in the PACE campaign. The two panels separate the five conceptual comparison fields for readability.}
\label{tab:prior}
\begin{tabular}{>{\raggedright\arraybackslash}p{2.1cm}>{\raggedright\arraybackslash}p{5.0cm}>{\raggedright\arraybackslash}p{7.1cm}}
\toprule
Method & Source-model supervision & Target information used \\
\midrule
MCDE~\cite{yue2026mcde} & Labeled benign and known-attack training; benign-mode fitting & Benign reference and threshold sets; trusted target benign recalibration in device transfer. \\
Kitsune~\cite{kitsune2018} & Unsupervised autoencoders; training traffic assumed benign & Initial local traffic for feature grouping and reconstruction learning. \\
D\"IoT~\cite{diot2019} & Benign device-type sequence learning across clients & Local device traffic and device type; threshold study uses benign and attack evaluation traffic. \\
OSDN~\cite{osdn2024} & Labeled source benign and attack categories & Unlabeled target traffic, potentially including unknown attacks; heterogeneous feature spaces. \\
PACE & Labeled benign and attack data at peers & Trusted benign fitting, reference, and calibration roles; pooled benign preprocessing. \\
\midrule
Method & Fusion or adaptation mechanism & Calibration or final decision \\
\midrule
MCDE & Supervised score plus distance to benign modes & Benign-tail alignment and branch allocation; separate final fused-score threshold. \\
Kitsune & Ensemble reconstruction errors feed an output autoencoder & Threshold on output reconstruction error; the paper discusses application-dependent threshold choices. \\
D\"IoT & Federated aggregation of device-type sequence models & Low packet-symbol probabilities identify anomalies; their fraction in a packet window triggers an alert. \\
OSDN & Heterogeneous projections and open-set domain alignment & Shared classifier trained on known classes and generated unknown examples. \\
PACE & Retained peer experts plus local anchor and expert bank & Separate peer, anchor, and bank calibration; final alert OR without aggregate recalibration. \\
\bottomrule
\end{tabular}
\end{table}


\subsection{Empirical calibration and statistical guarantees}
Rank-based score calibration has a substantial statistical foundation. Split conformal prediction uses held-out calibration scores and a finite-sample quantile correction to obtain marginal coverage under exchangeability; the probability averages over calibration and test sampling~\cite{angelopoulos2023conformal}. This is different from merely ensuring that a prescribed fraction of the observed calibration scores exceeds an empirical threshold. PACE uses the latter type of empirical constraint with a maximum over block-specific quantiles and data-dependent union scaling. Its threshold index is not the standard split-conformal correction, and its extrapolated evidence is not itself a conformal p-value.

Bates et al.~\cite{bates2023outliers} distinguish marginal conformal outlier p-values from calibration-conditional validity. Shared calibration observations induce dependence between test p-values, which matters for downstream multiple testing. Their conditional constructions require explicit statistical adjustments. Neither calling an equal-share detector union ``Bonferroni'' nor using a smoothed empirical tail establishes those properties for PACE. False discovery rate across tested observations also differs from a detector's benign false-alert probability; the quantities should not be interchanged. Formally valid p-value combination requires valid marginal inputs, even when a combination rule accommodates arbitrary dependence~\cite{vovk2020averaging}. The union bound itself needs no independence assumption, but empirical threshold calibration does not establish the population tail bounds needed to apply it at a claimed test-time level.

Barber et al.~\cite{barber2023nonexchangeable} extend conformal prediction with weighted quantiles and bounds on the loss of coverage under departures from exchangeability. Their results connect validity to the difference between the relevant data distributions and use a specifically constructed prediction rule. PACE does not implement that procedure. Its temporal partition and worst-block calibration do not by themselves import protection against drift. These distinctions motivate reporting realized false-alert distributions and limiting conclusions to the tested setting.

A recent CAN preprint by Xiang et al.~\cite{xiang2026can} separates benign target adaptation from later benign calibration and applies empirical ranks to temporal block maxima. Its block length reflects estimated serial dependence and available rank resolution; unresolved dependence limits the validity claim. PACE instead takes the largest empirical quantile across calibration blocks and applies the resulting threshold to individual records. These operations use different decision units. The preprint reinforces that benign-only adaptation and block-based empirical calibration are existing alternatives to supervised target fitting; it does not supply a guarantee for PACE's peer unions.
\subsection{Established local anomaly models}
PACE uses established anomaly-scoring mechanisms rather than a new base detector. Isolation forest scores the ease of isolating an observation in random partition trees~\cite{liu2008isolation}; local outlier factor compares neighborhood density with that of nearby observations~\cite{breunig2000lof}. These models encode different notions of deviation, but their diversity alone does not establish that a combined detector will improve detection at a fixed false-alert rate. That question requires evaluating the complete scoring and thresholding rule.

The seven-member comparator also includes reconstruction error from an autoencoder. Pang et al.~\cite{pang2021review} distinguish reconstruction-based learning from objectives designed directly for anomaly scoring: good compression of dominant regularities need not preserve the information needed to separate anomalies. This motivates treating reconstruction as one source of evidence, but it does not explain PACE's measured gains. PACE's six-member anchor omits the autoencoder while changing local calibration and adding a bank, so the primary result cannot isolate an autoencoder effect.

\subsection{Benign-only monitoring, transfer, and personalization}
Benign-only target modeling is established in IoT security. N-BaIoT fits device-local autoencoders to benign traffic and uses separate benign observations to choose detection settings~\cite{nbaiot2018}. Its nine-device Mirai and BASHLITE evaluation is relevant precedent, but its packet and window decisions cannot be ranked directly against the present target-macro evaluation. D\"IoT instead learns device-type traffic models collaboratively through federated aggregation~\cite{diot2019}. PACE assumes supervised source models trained on labeled peer attacks and retains those models separately; it does not perform iterative federated optimization.

FedMSE~\cite{nguyen2025fedmse} learns benign-only autoencoder representations with centroid anomaly scoring and averages model parameters using inverse reconstruction error on a shared processed benign set. Its collaborative learning differs from retaining PACE's separately trained supervised peers. The distinction concerns the learned representation, aggregation rule, and information access.

Recent personalization methods also learn client-specific model parameters. The federated hypernetwork IDS of Abhijit et al.~\cite{hyperids2025} updates a shared hypernetwork and client representations under labeled classification objectives. PACE personalizes evidence and thresholds from target benign observations while freezing the peer classifiers. The methods therefore personalize different parts of the detector under different information regimes. The centralized pooled preprocessing used in PACE's benchmark further limits any literal model-only deployment interpretation.

Open Set Dandelion Network (OSDN) learns from labeled source data and unlabeled target data through heterogeneous projections and domain alignment~\cite{osdn2024}. Its target data can include unknown attacks. PACE instead requires trusted benign target roles and a compatible feature schema for the supplied peer models. The comparison distinguishes two attack-label-free target settings: adaptation on mixed unlabeled traffic with learned cross-schema alignment, and adaptation on trusted benign traffic with frozen compatible peer models.

Calibration also has different meanings across IDS studies. Talpini et al.~\cite{talpini2025} federate parameters of sigmoid probability calibration using labeled local calibration data and a classification loss. Their goal concerns the reliability of class probabilities. PACE calibrates alert thresholds against benign empirical tails. A better-calibrated class probability need not meet a chosen benign false-alert rate, and an empirical tail threshold need not yield calibrated posterior probabilities.

\subsection{Fusion and detector-response representations}
Zimek et al.~\cite{zimek2014} discuss detector diversity, score normalization, and the sensitivity of maximum fusion to an individual overestimating detector. PACE's peer-wise thresholds address a particular consequence of these issues when transferred peer scores are calibrated on a new target. The shared union multiplier still couples the peers: separate thresholds do not make the final decision immune to another expert's behavior.

Modeling joint detector responses also has direct network-security precedent. Kitsune passes the normalized reconstruction errors of an autoencoder ensemble to another autoencoder, allowing relationships among first-layer errors to influence anomaly detection~\cite{kitsune2018}. PACE's response expert instead fits an isolation forest to target-relative evidence from foreign supervised peers. The choice of response coordinates and estimator differs; the general idea of anomaly detection in detector-output space is not new. ViM combines classifier logits with feature-space residual information for out-of-distribution detection~\cite{vim2022}, providing a further example of complementary output and representation evidence. Its vision evaluation and geometry differ from PACE's peer-response forest and do not constitute a direct IDS baseline.

\subsection{Transfer evaluation and operating points}
Recent cross-protocol evidence reinforces the need to separate ranking from threshold transfer. Jin~\cite{jin2026crossprotocol} evaluates IP-to-BLE adaptation using unlabeled target data, contrasting random-window evaluation with a capture-wise hold-out and auditing source-selected thresholds. Improvements depend on the split, while transferred thresholds can produce severe false-alert failures. That setting differs from PACE's trusted target-benign threshold adaptation and within-stratum feature compatibility. Nevertheless, it illustrates why a pooled detection gain or an unlabeled adaptation procedure does not establish safe operation on a new capture. PACE's false-alert distributions and preprocessing dependencies must be evaluated on their own evidence.



Ha and Kim~\cite{ha2026targetfpr} calibrate benign-validation thresholds at specified false-positive rates and transfer them unchanged to held-out attacks in host-session graph detection. They report realized false-positive rates relative to the requested level and examine failure in the tails of that distribution. Their models use labeled benign and seen-attack training examples; PACE instead studies a benign-only target's reuse of separate supervised peers and calibration of their alert union. The shared principle of assessing a fixed threshold on unseen attacks does not make the studies' aggregation units or numerical gains directly comparable.